\section{Introduction}

A benchmark score can be perfectly computed and still fail to support the interpretation attached
to it. Construct underrepresentation, shortcuts, contamination, scoring artifacts, dependent
items, saturation, and weak discrimination affect different inferences and should not be collapsed
into one health score. This distinction is familiar in validity theory
\citep{messick1995validity,aera2014standards} but remains unevenly enforced in contemporary model
evaluation.

ValidEval asks a prior question: who validates the diagnostics used to audit a benchmark? A warning
that correlates with model disagreement is not automatically an error detector. A response matrix
with many models is not automatically a valid IRT sample. A large rank range is not automatically a
material decision change. Our central claim is methodological: diagnostic outputs should be linked
to explicit estimands, uncertainty, validation evidence, and blocked wording.

We separate a historical imported study (Study H) from a controlled exact-checkpoint study (Study
C). The current paper scaffold reports only artifact-backed Study H results. Every unexecuted study
retains an exact result requirement rather than a plausible number.
